% LIPIcs skeleton for ESEM 2026 Emerging Results submission.
% Double-anonymous. Do not include author names, affiliations, or grants.
\documentclass[a4paper,UKenglish,cleveref,autoref,thm-restate]{lipics-v2021}

\bibliographystyle{plainurl}

\title{When Does Adaptive Query Execution Pick the Wrong Join?\\
       An Empirical Study of Strategy Selection in Apache Spark}
\titlerunning{When Does AQE Pick the Wrong Join?}

% Anonymous author block.
\author{Anonymous}{Anonymous}{}{}{}
\authorrunning{Anonymous}
\Copyright{Anonymous}

\ccsdesc[500]{Information systems~Query optimization}
\ccsdesc[300]{Software and its engineering~Software performance}

\keywords{Apache Spark, adaptive query execution, joins, empirical evaluation}

\category{Emerging Results}
\relatedversion{}
\supplement{Replication package at \url{<ANONYMIZED_URL>}}

\hideLIPIcs

\begin{document}
\maketitle

\begin{abstract}
Apache Spark's Adaptive Query Execution (AQE) re-decides physical join
strategies at runtime using statistics gathered after intermediate shuffles.
We empirically measure how often this decision matches an oracle obtained by
forcing each viable strategy and characterise the workload features that
predict mismatches. On TPC-DS SF=1000 and TPC-H SF=1000, AQE's choice differs
from the fastest viable strategy on a non-trivial fraction of join sites.
Mismatches are concentrated in plans where the estimated and actual build-side
sizes diverge and where the post-shuffle partition size exceeds the L3 cache.
We demonstrate that a configuration-level adjustment recovers part of the gap
without code changes, and we outline a longer-term research agenda toward
cache-aware adaptive planning.
\end{abstract}

\section{Introduction}
% Motivation: AQE re-plans at runtime. Open question: does it pick the best?
% Contribution bullets:
%  - empirical decision-accuracy measurement on TPC-DS and TPC-H
%  - feature-level characterisation of failures
%  - config-level mitigation as initial evidence
% Scope: emerging results; full system fix is future work.

\section{Background}
\subsection{Join physical operators in Spark}
% Brief: BHJ, SHJ, SMJ. Memory and network costs.
\subsection{Adaptive Query Execution}
% Brief: DynamicJoinSelection, OptimizeSkewedJoin, CoalesceShufflePartitions.

\section{Research Questions}
\label{sec:rq}
% RQ1, RQ2, RQ3 verbatim from docs/research-questions.md.

\section{Method}
\label{sec:method}
% Cluster, workloads, planner modes, repetitions, statistical analysis.
% Forward-reference docs/methodology.md.

\section{Results}
\subsection{RQ1: AQE decision accuracy}
\subsection{RQ2: Predictors of mismatch}
\subsection{RQ3: Configuration-level mitigation}

\section{Threats to Validity}
% Construct / internal / external / conclusion. Forward-reference
% docs/threats-to-validity.md condensed.

\section{Related Work}
% Cache-aware joins in DBs (third-person): Shatdal et al., Kim et al., Balkesen
% et al., Schuh et al. Plus AQE / Photon / Velox style runtime adaptation.

\section{Discussion and Future Work}
% Failure mode taxonomy implies system work: cache-aware partition sizing,
% in-task radix re-partitioning, spillable broadcast hash relation.

\section*{Data Availability}
All scripts, configurations, query lists, and analysis notebooks are available
as an anonymised replication package at \url{<ANONYMIZED_URL>}. Raw Spark
event logs and processed CSV summaries are included.

\bibliography{references}

\end{document}
